本文面向硬禾科技 eZ-PLM 平台，设计并实现了一套融合 LLM 与工程规则推理的智能元器件选型与风险评估 Agent 系统，实现了从自然语言需求到 BOM 选型清单、风险评估报告和拓扑分析报告的端到端自动化流程。系统采用 Pipeline 与 ReAct Agent 双模架构，兼顾批量效率与多轮交互推理能力。

在数据集成方面，系统完成了 eZ-PLM API 全链路接入，构建了基于制造商型号前缀的分组查询与 HMAC-SHA256 签名认证机制，通过多前缀检索、五维约束过滤和 MPN 电压推断实现了对在产器件的准确召回。

在需求理解方面，提出四级容错需求解析链，将 LLM 语义解析与规则校验结合，28 条评测用例实现 100\% 解析通过率，验证了该方法在工程场景中的鲁棒性。

在智能评分方面，构建了参考设计知识增强的双模自适应评分框架，参考设计可用时引入 LLM 软评分维度，数据缺失时自动回退至纯规则模式，兼顾评分稳定性与知识增强能力。

在风险评估方面，提出规则验证与 LLM 叙述分离的双层架构，从机制上避免语言模型不确定性对风险结论的影响。

在架构层面，建立了以 \texttt{PartIR}、\texttt{ScoreBreakdown}、\texttt{TopologyIR} 和 \texttt{RiskIR} 为核心的标准化 IR 体系，同步生成结构化 JSON 与标准化工程文档，支持 eZ-PLM 平台数据互通。同时，基于 ChromaDB 构建工程知识库，结合 LangChain ReAct Agent 实现知识检索与替代分析等能力的自然语言统一调用。

实验结果表明，本方法有效提升了元器件选型与风险评估的自动化程度和结果可解释性，为 LLM 在电子设计自动化与产品生命周期管理场景中的应用提供了工程可行方案。
